\chapter{Conductors and Capacitors} \label{chap:conductors-and-capacitors}

\section{Conductors}

\section{Capacitors} \label{sec:capacitors}

\textbf{Capcitors} are devices that store electric charge.
Unlike batteries, they are purely electrical rather than electrochemical.
Physically,
they are composed of two conductors separated by a distance,
which have opposite charges of equal magnitude when charged.

How well a capacitor stores charge is its \textbf{capacitance} (\(C\)),
with SI units of farads \unit{\farad}.
The capacitance is defined as
\begin{equation}
  C = \frac{q}{V},
\end{equation}
where \(q\) is the magnitude of charge on each conductor
and \(V\) is the voltage between the two conductors.

\subsection{Parallel Plate Capcitors}

\textbf{Parallel plate capacitors} are the most common type of capacitor.
They are composed of two parallel plates,
and their capacitance \(C\) is given by
\begin{equation}
  C = \frac{\epsilon_\mathrm{d} A}{d},
\end{equation}
where \(\epsilon_\mathrm{d}\) is the permittivity of the \textbf{dielectric},
\(A\) is the area of each plate,
and \(d\) is the distance between the plates.
